\subsection{Integration with respect to the image of a positive measure}

Let \(\pi\) be a \(\mu\) -proper mapping of T into X, and let \(\nu = \pi(\mu)\) . Applying the results of §4, one obtains the following statements:

PROPOSITION 2. --- For every numerical function \(f \geqslant 0\) defined on X,

\[
\int^ {\bullet} f (x) d \nu (x) = \int^ {\bullet} f (\pi (t)) d \mu (t).\tag{2}
\]

This follows from Th. 1 of §4, No.~2.

COROLLARY 1. --- For every subset A of X,

\[
\nu^ {\bullet} (\mathrm{A}) = \mu^ {\bullet} \left(\overline {{\pi}} ^ {- 1} (\mathrm{A})\right).
\]

COROLLARY 2. --- For a subset A of X to be locally negligible for \(\nu\) , it is necessary and sufficient that \(\overline{\pi}^{-1}(\mathrm{A})\) be locally negligible for \(\mu\) .

COROLLARY 3. --- If the measure \(\mu\) is concentrated on a set M, then \(\pi(\mu)\) is concentrated on \(\pi(M)\) .

For, if \(\mathrm{N} = \mathrm{X} - \pi(\mathrm{M})\) then \(\overline{\pi}^1(\mathrm{N})\) does not intersect \(\mathrm{M}\) , therefore is locally \(\mu\) -negligible, consequently (Cor. 2) \(\mathrm{N}\) is locally \(\nu\) -negligible.

COROLLARY 4. --- Let S be the support of \(\mu\) . If \(\pi\) is continuous, then the support of \(\pi(\mu)\) is \(\overline{\pi(S)}\) .

For, it follows from Cor. 3 that \(\pi(\mu)\) is concentrated on \(\pi(S)\) , therefore if \(S'\) is the support of \(\pi(\mu)\) then \(S' \subset \overline{\pi(S)}\) . On the other hand, \(\overline{\pi}(X - S')\) is a locally \(\mu\) -negligible open set (Cor. 2), hence is \(\mu\) -negligible (Ch. IV, §5, No.~2, Cor. 2 of Prop. 5). Therefore \(\overline{\pi}(X - S') \subset T - S\) , consequently \(\pi(S) \subset S'\) , which proves the corollary.

PROPOSITION 3. --- For a mapping f of X into a topological space G to be ν-measurable, it is necessary and sufficient that \(f \circ \pi\) be μ-measurable.

This is an immediate consequence of Prop. 3 of §4, No.~3.

COROLLARY. --- For a subset A of X to be \(\nu\) -measurable, is is necessary and sufficient that \(\overline{\pi}^{1}(A)\) be \(\mu\) -measurable.

\pandocbounded{\includegraphics[keepaspectratio]{images/part002_e5542695.jpg}}

However, the image under \(\pi\) of a \(\mu\) -measurable subset M of T is not necessarily \(\nu\) -measurable, even if \(\pi\) is continuous and M is \(\mu\) -negligible (Exer. 7 and §8, Exer. 1).

THEOREM 1. --- Let f be a function defined on X with values in \(\overline{R}\) or in a Banach space F. For f to be essentially \(\nu\) -integrable, it is necessary and sufficient that \(f \circ \pi\) be essentially \(\mu\) -integrable, in which case

\[
\int \mathbf {f} (x) d \nu (x) = \int \mathbf {f} (\pi (t)) d \mu (t).\tag{3}
\]

Suppose, moreover, that \(\pi\) is continuous and proper. Then, for \(\mathbf{f}\) to be \(\nu\) -integrable, it is necessary and sufficient that \(\mathbf{f} \circ \pi\) be \(\mu\) -integrable.

It suffices to apply Th. 2 of §4, No.~4.

COROLLARY. --- For a subset A of X to be essentially \(\nu\) -integrable, it is necessary and sufficient that \(\overline{\pi}^{1}(A)\) be essentially \(\mu\) -integrable, in which case \(\nu(A)=\mu\left(\overline{\pi}^{1}(A)\right)\) .

In particular, for every compact set \(K \subset X\) , \(\nu(K) = \mu\left(\overline{\pi}(K)\right)\) . It follows from this and Cor. 3 of Prop. 2 that if \(\mu\) is atomic ( \(\S5\) , No.~10) then so is \(\pi(\mu) = \nu\) . For, let M be the set of \(t \in T\) such that \(\mu(\{t\}) \neq 0\) ; since \(\mu\) is carried by M, \(\nu\) is carried by \(\pi(M)\) ; moreover, for every \(x \in \pi(M)\) we have \(\nu(\{x\}) = \mu\left(\overline{\pi}(x)\right) > 0\) , since \(\overline{\pi}(x)\) contains at least one point of M. Therefore \(\nu\) is atomic ( \(\S5\) , No.~10, Prop. 15).

